% ==============================================================================
% 03_dataset.tex
% Phụ trách: Thành viên 1 (Cơ sở nghiên cứu, phương pháp và dữ liệu)
% ==============================================================================
% TODO [Thành viên 1]:
% - Mô tả corpus 384 tài liệu Angular Markdown, cấu trúc tài liệu và tiền xử lý.
% - Quá trình dùng ChatGPT hỗ trợ tạo QA ban đầu.
% - Quá trình 3 kỹ sư phần mềm review tuần tự (sequential expert review).
% - Cấu trúc mỗi QA record & NGUYÊN TẮC QUAN TRỌNG: không gắn ground truth với chunk ID
%   (để đảm bảo tính trung lập khi đánh giá các chiến lược chunking khác nhau).
% - Bảng thống kê chi tiết: 140 câu hỏi, 233 evidence items, 60 easy, 40 medium, 40 hard.
% - Phân bố phạm vi: single-section, cross-section, cross-document.
% - Giải thích tiêu chí thiết kế độ khó và các nhóm câu hỏi trong benchmark.
% ==============================================================================

% 03_dataset.tex
% Phụ trách: Thành viên 1
% ==============================================================================

\section{Corpus and QA Benchmark}
\label{sec:dataset}

\subsection{Technical Document Corpus}

Corpus gồm 384 tài liệu kỹ thuật Angular ở định dạng Markdown. Các tài liệu có cấu trúc phân cấp thông qua heading và chứa nhiều loại block như paragraph, list, table, code block và callout. Đặc điểm này cho phép đánh giá ảnh hưởng của việc bảo toàn cấu trúc tài liệu trong quá trình chunking.

Trong bước tiền xử lý, mỗi tài liệu được phân tích thành chuỗi block theo thứ tự xuất hiện trong nguồn. Đối với mỗi block, hệ thống lưu loại nội dung, văn bản gốc và đường dẫn heading tương ứng. Nội dung không được tóm tắt hoặc viết lại nhằm bảo toàn bằng chứng nguồn cho quá trình đánh giá.

\subsection{QA Benchmark Construction}

Do chưa có benchmark phù hợp trực tiếp với corpus và mục tiêu nghiên cứu, nhóm xây dựng một bộ câu hỏi--trả lời có gắn bằng chứng nguồn. ChatGPT được sử dụng để hỗ trợ tạo câu hỏi ứng viên, đáp án tham chiếu và vùng bằng chứng ban đầu. Các đầu ra này chỉ được xem là bản nháp và không được đưa trực tiếp vào benchmark cuối cùng.

Mỗi mẫu sau đó trải qua \textit{sequential expert review} bởi ba kỹ sư phần mềm có từ bốn đến năm năm kinh nghiệm. Các chuyên gia lần lượt kiểm tra và chỉnh sửa câu hỏi, đáp án, loại câu hỏi, kiểu suy luận, độ khó và bằng chứng nguồn. Một mẫu chỉ được chấp nhận khi có thể trả lời đầy đủ từ corpus và không phụ thuộc vào kiến thức bên ngoài.

\subsection{Data Structure and Chunk-Neutral Ground Truth}

Mỗi QA record gồm các trường chính sau:

\begin{itemize}
    \item \texttt{question\_id}: mã định danh của câu hỏi;
    \item \texttt{question}: nội dung câu hỏi;
    \item \texttt{reference\_answer}: đáp án tham chiếu;
    \item \texttt{difficulty}: mức độ easy, medium hoặc hard;
    \item \texttt{question\_type}: loại câu hỏi;
    \item \texttt{reasoning\_type}: kiểu suy luận cần thiết;
    \item \texttt{evidence}: danh sách các evidence item.
\end{itemize}

Mỗi evidence item lưu mã bằng chứng, mã tài liệu, đường dẫn section và các câu chứa bằng chứng. Ground truth được liên kết với nội dung tài liệu gốc thay vì chunk được tạo bởi một phương pháp cụ thể. Benchmark không lưu \texttt{chunk\_id}, \texttt{relevant\_chunk\_id} hoặc nhãn phụ thuộc vào chiến lược chunking. Nguyên tắc này cho phép cùng một ground truth được sử dụng công bằng cho fixed-size, structure-aware và prompt-based chunking.

\subsection{Benchmark Statistics}

Benchmark cuối cùng gồm 140 câu hỏi và 233 evidence items. Phân bố theo độ khó được trình bày trong Bảng~\ref{tab:difficulty-distribution}.

\begin{table}[t]
\centering
\small
\begin{tabular}{lrr}
\toprule
\textbf{Độ khó} & \textbf{Số câu} & \textbf{Tỷ lệ} \\
\midrule
Easy   & 60  & 42.9\% \\
Medium & 40  & 28.6\% \\
Hard   & 40  & 28.6\% \\
\midrule
Tổng   & 140 & 100.0\% \\
\bottomrule
\end{tabular}
\caption{Phân bố câu hỏi theo độ khó.}
\label{tab:difficulty-distribution}
\end{table}

Các câu hỏi còn được phân loại theo phạm vi bằng chứng. \textit{Single-section} sử dụng bằng chứng trong một section; \textit{cross-section} kết hợp nhiều section trong cùng tài liệu; và \textit{cross-document} yêu cầu bằng chứng từ nhiều tài liệu. Phân bố tương ứng được trình bày trong Bảng~\ref{tab:evidence-scope}.

\begin{table}[t]
\centering
\small
\begin{tabular}{lr}
\toprule
\textbf{Phạm vi bằng chứng} & \textbf{Số câu} \\
\midrule
Single-section & 50 \\
Cross-section  & 60 \\
Cross-document & 30 \\
\midrule
Tổng           & 140 \\
\bottomrule
\end{tabular}
\caption{Phân bố câu hỏi theo phạm vi bằng chứng.}
\label{tab:evidence-scope}
\end{table}

\subsection{Difficulty Design}

Độ khó được xác định dựa trên số vùng bằng chứng, độ phức tạp suy luận, mức độ phụ thuộc giữa các phần tài liệu và phạm vi thông tin cần bao phủ. Tất cả câu hỏi đều phải có đủ bằng chứng trong corpus; nhãn hard không biểu thị câu hỏi mơ hồ hoặc thiếu thông tin.

Câu hỏi easy chủ yếu thuộc nhóm fact hoặc definition và dựa trên một vùng bằng chứng cục bộ. Câu hỏi medium yêu cầu kết hợp hai evidence items, thường thuộc các nhóm comparison, procedure, cause--effect hoặc multi-fact synthesis. Trong 40 câu medium, 27 câu có hai evidence items trong cùng một section và 13 câu sử dụng bằng chứng từ nhiều section. Câu hỏi hard kết hợp từ hai đến ba vùng bằng chứng và có thể yêu cầu tổng hợp, phân tích trade-off, theo dõi thay đổi trạng thái hoặc đối chiếu thông tin phân tán.